\documentclass[aps,prd,reprint,nofootinbib,superscriptaddress,longbibliography]{revtex4-2}

\usepackage{amsmath,amssymb,mathtools,graphicx,bm,booktabs}
\usepackage[colorlinks=true,allcolors=blue]{hyperref}

\newcommand{\std}{\mathrm{std}}
\newcommand{\ESS}{\mathrm{ESS}}
\newcommand{\maybegraphic}[3][0.95\linewidth]{%
  \IfFileExists{#2}{\includegraphics[width=#1]{#2}}{%
    \fbox{\parbox[c][2.0in][c]{0.88\linewidth}{\centering \textbf{Figure placeholder}\\[0.5em]#3}}%
  }%
}

\begin{document}

\title{Learned Wilsonian Hierarchies and Renormalized Trajectories in Two-Dimensional Lattice $\phi^4$ Theory}

\author{Kostas Orginos}

\begin{abstract}
We propose a Wilsonian use of hierarchical normalizing flows for lattice field
theory that goes beyond exact reweighting. The starting point is an invertible
multiscale transport map for two-dimensional lattice $\phi^4$ theory in which
each $2\times2$ blocking step is decomposed into one coarse scalar field and
three local fluctuation modes, and only the fluctuations are trivialized while
the coarse field is held fixed. This gives the learned model a hierarchical
factorization in which every RG level defines an exact blocked marginal of the
learned action, even when the explicit coarse effective action is never
parametrized. We argue that this provides an exact sampler for the blocked
theories of the learned model and therefore an operational route to learned
renormalization-group trajectories in measure space. The central hypothesis is
that, after training near a target point $(g_2,g_4)$, the resulting blocked
hierarchy flows toward a renormalized trajectory associated with a nearby point
in coupling space, even when fine-level reweighting becomes impractical because
the residual mismatch is extensive. This paper sets up the formalism, the
observable-space RG program, and the numerical strategy needed to test this
claim in two-dimensional $\phi^4$ theory.
\end{abstract}

\maketitle

\section{Introduction}

Normalizing flows in lattice field theory have been used primarily as exact or
approximately exact sampling tools \cite{Dinh2017RealNVP,Albergo2019FlowMCMC,Kanwar2020Equivariant,Albergo2021FlowScaling}.
The usual viewpoint is microscopic: train a flow to match a target measure,
sample from the flow, and correct any residual mismatch by reweighting. This is
powerful, but it eventually runs into an extensive problem. If the residual
action difference is a sum of small local mismatches, then
\begin{equation}
\Delta S(\phi)=S(\phi)-S_\theta(\phi)
\end{equation}
has a width that typically grows like $\sqrt{V}$, and the reweighting factors
become broad even when the learned action is still close to the target action in
an intensive sense.

This suggests a different use of the learned hierarchy. If the transport map is
constructed in a Wilsonian way, with the coarse field held fixed and only the
block fluctuations trivialized, then the model naturally defines a family of
blocked measures at all RG levels. These blocked measures are exact marginals of
the learned action. Therefore one does not need to parametrize the coarse
effective action in order to sample it. One can instead study the blocked
theories directly in the space of observables or, after an additional
projection, in a truncated coupling space.

The present manuscript is the conceptual companion to the shorter methods paper.
There the emphasis is on the transport map itself, learned Gaussian priors, and
the scaling of the residual mismatch. Here the emphasis is different:
\begin{enumerate}
\item the learned hierarchy as an exact blocked sampler for the learned action,
\item RG trajectories in observable space without explicit coarse-action
parametrization,
\item the relation between approximate matching of the target action and
universality,
\item the possibility of replacing ``simulate the target action exactly'' by
``simulate a nearby representative of the same renormalized trajectory.''
\end{enumerate}

\section{Exact Blocked Hierarchy of the Learned Action}
\label{sec:exact_blocked_hierarchy}

Let $q_\theta(\phi^{(0)})$ be the learned fine-lattice measure defined by the
hierarchical transport map. At each $2\times2$ block we rewrite the field in
terms of a coarse scalar and three local fluctuations,
\begin{equation}
\phi^{(k)} \longleftrightarrow \left(\phi^{(k+1)},\eta^{(k)}\right).
\end{equation}
Because the transport map is built recursively and remains exactly invertible,
the model measure factorizes exactly as
\begin{equation}
q_\theta(\phi^{(0)})
=
q_\theta(\phi^{(K-1)})
\prod_{k=0}^{K-2}
q_\theta\!\left(\eta^{(k)} \mid \phi^{(k+1)}\right).
\label{eq:q_factorization}
\end{equation}

Equation~\eqref{eq:q_factorization} has an immediate consequence: for each RG
level $k$, the blocked measure
\begin{equation}
q_\theta^{(k)}(\phi^{(k)})
\end{equation}
is an exact marginal of the learned action. Therefore:
\begin{itemize}
\item it is sampled exactly by ancestral generation through the learned
hierarchy,
\item it defines an exact blocked effective action
\begin{equation}
S_\theta^{(k)}(\phi^{(k)})=-\log q_\theta^{(k)}(\phi^{(k)}),
\end{equation}
even if that action is never written in closed form,
\item one can study the RG flow of the learned theory directly without fitting a
separate coarse action.
\end{itemize}

This point is conceptually important. In the standard variational RG problem,
one often introduces an ansatz for the blocked action and then fits its
parameters. Here the blocked measures exist exactly for the learned action before
any such projection is performed.

\begin{figure}[t]
  \centering
  \maybegraphic{figures/blocked_sampler_hierarchy.pdf}{Planned figure: exact
  blocked-sampler view of the learned hierarchy, showing the fine model and all
  coarse marginals generated by ancestral sampling.}
  \caption{The learned hierarchy as an exact blocked sampler for the learned
  action. Even without an explicit closed-form coarse action, each blocked level
  defines an exact marginal of the learned measure and can therefore be sampled
  directly.}
  \label{fig:blocked_sampler}
\end{figure}

\section{Target Theory Versus Learned Theory}
\label{sec:target_vs_learned}

The construction above is exact for the learned measure $q_\theta$, not
necessarily for the target measure $p \propto e^{-S}$. If the learned action
were exact, $q_\theta=p$ and the blocked hierarchy would coincide with the true
Wilsonian RG flow of the target theory. In practice, one has instead a residual
mismatch
\begin{equation}
\Delta S(\phi)=S(\phi)-S_\theta(\phi).
\end{equation}

The key observation from the current numerical evidence is that this mismatch can
be small in an intensive sense even when fine-level reweighting is already
difficult. If $\Delta S$ is approximately a sum of local terms,
\begin{equation}
\Delta S(\phi) \approx \sum_x \delta s_x(\phi),
\end{equation}
then
\begin{equation}
\mathrm{Var}(\Delta S)\propto V,
\qquad
\std(\Delta S)\propto \sqrt{V}.
\end{equation}
In this regime, exact reweighting eventually becomes costly for purely
extensive reasons, while the learned action may still be close to the target in
the space of local theories.

This distinction motivates the present program. Instead of treating a broad
reweighting distribution as proof that the learned theory is unusable, one asks a
different question: is the learned blocked hierarchy flowing toward the same
renormalized trajectory as the target theory?

\section{Observable-Space RG Trajectories}
\label{sec:observable_space}

The first version of this program should avoid any unnecessary projection onto a
specific coupling basis. The natural first object is therefore the RG trajectory
in the space of observables.

For each bare point $(g_2,g_4)$:
\begin{enumerate}
\item train a Wilsonian transport map near that target point,
\item sample from the blocked marginals $q_\theta^{(k)}$ at every RG level,
\item measure a vector of coarse observables at each level.
\end{enumerate}

Natural observables include:
\begin{itemize}
\item the second-moment correlation length ratio $\xi/L$,
\item the Binder cumulant
\begin{equation}
U_4 = \frac{\langle M^4\rangle}{\langle M^2\rangle^2},
\qquad
M = \frac{1}{V}\sum_x \phi_x,
\end{equation}
\item susceptibilities,
\item blocked two-point functions at small momenta,
\item connected four-point cumulants.
\end{itemize}

This produces a discrete trajectory
\begin{equation}
\mathcal O^{(0)} \to \mathcal O^{(1)} \to \cdots \to \mathcal O^{(K-1)}
\end{equation}
for the learned theory. If two different fine-level target points yield blocked
trajectories that merge after several RG levels, then one has direct evidence
that the learned hierarchy is following a common renormalized trajectory in
observable space.

\section{Projection to Coupling Space}
\label{sec:coupling_space}

Once the observable-space trajectories are understood, one may project the
blocked theory onto a truncated coupling basis,
\begin{equation}
S_{\mathrm{ansatz}}^{(k)}(\varphi)
=
\sum_i g_i^{(k)} O_i(\varphi),
\end{equation}
for example using the operator set
\begin{equation}
\phi^2,\quad \phi^4,\quad \phi^6,\quad
\sum_{x,\mu} \phi_x \phi_{x+\hat\mu},
\quad \ldots
\end{equation}
This projection is useful, but it should be viewed as a second step. The
blocked measures of the learned theory exist before any truncation and are exact
within the learned model.

In this sense the coupling-space description is diagnostic rather than
foundational: it provides coordinates on the trajectory, but it is not required
for the existence of the trajectory itself.

\section{Universality Interpretation}
\label{sec:universality}

For two-dimensional $\phi^4$ theory with exact $Z_2$ symmetry, the critical line
in the bare $(g_2,g_4)$ plane flows into the Ising universality class in the
infrared. The precise microscopic representative changes with the bare couplings,
but long-distance observables on the critical surface share the same universal
behavior.

This suggests a nonstandard but physically natural objective for learned
transport maps. One need not insist that the learned action reproduce a chosen
microscopic target action with arbitrarily good fine-level reweighting on very
large volumes. It may be enough that the learned blocked hierarchy lies on, or
very near, the same renormalized trajectory. Then the learned theory serves as a
practical representative of the same universality class.

From this point of view, the learned hierarchy functions as a kind of
\emph{quantum perfect action} construction: not because it reproduces the exact
continuum fixed-point action in closed form, but because it provides exact
blocked samplers for a learned theory that may lie very close to the relevant
renormalized trajectory.

\section{Numerical Program}
\label{sec:numerical_program}

The numerical program for this paper has four stages.

\subsection{Stage 1: volume scaling}

Train the current Wilsonian transport map at
\begin{equation}
L=16,\ 32,\ 64,\ 128
\end{equation}
for a fixed canonical target point, and study:
\begin{equation}
\std(\Delta S),\qquad \std(\Delta S)/L,\qquad \ESS.
\end{equation}
The goal is to determine whether the residual mismatch remains intensive as the
volume grows.

\subsection{Stage 2: blocked observables}

At each RG level of the learned hierarchy, measure:
\begin{itemize}
\item $\xi/L$,
\item $U_4$,
\item low-momentum two-point observables,
\item connected four-point observables.
\end{itemize}
This gives the RG trajectory in observable space.

\subsection{Stage 3: target-point scan}

Repeat the construction for several bare points $(g_2,g_4)$, with particular
focus on scans in $m^2$ at fixed $\lambda$. The goal is to determine whether the
blocked trajectories merge and how the learned hierarchy maps the critical
surface.

\subsection{Stage 4: optional coupling projection}

Project the blocked theory onto a truncated local operator basis to extract a
coarse-grained flow in coupling space. This step is useful for interpretation
but is not required for the core claim.

\section{Figures for the Full Paper}
\label{sec:figures}

The natural figure set for this paper is:
\begin{enumerate}
\item schematic of the learned blocked-sampler hierarchy,
\item scaling of $\std(\Delta S)$ and $\std(\Delta S)/L$ with volume,
\item blocked observables versus RG depth for several bare points,
\item merging of observable-space trajectories,
\item optional projection of the same trajectories into a truncated
$(g_2,g_4,\ldots)$ space.
\end{enumerate}

\begin{figure}[t]
  \centering
  \maybegraphic{figures/rg_trajectory_observable_space.pdf}{Planned figure:
  observable-space RG trajectories for several target points, showing possible
  convergence toward a common renormalized trajectory.}
  \caption{Target figure for the central conceptual claim of the paper: the
  blocked hierarchies generated by learned Wilsonian transport maps may flow
  toward a common renormalized trajectory in observable space even when exact
  fine-level reweighting is already difficult.}
  \label{fig:observable_rg}
\end{figure}

\section{Scope of the Present Scaffold}
\label{sec:scope}

This manuscript is currently a scaffold rather than a finished paper. Its role
is to fix the conceptual structure of the RG/universality project:
\begin{itemize}
\item the exact blocked hierarchy of the learned action,
\item the distinction between target and learned measures,
\item observable-space trajectories as the first-class object,
\item coupling-space projections as a later diagnostic step,
\item universality as the long-distance criterion of success.
\end{itemize}

The actual paper should only be completed once the large-volume runs and the
observable-space RG scans are in hand.

\section{Outlook}
\label{sec:outlook}

If this program succeeds, it will change the role of normalizing flows in
lattice field theory. The point will no longer be only to construct a transport
map that reweights efficiently at the fine level. The point will be to construct
a learned Wilsonian hierarchy that gives access to renormalized trajectories and
coarse effective theories operationally, without an explicit coarse-action
parametrization. That would make the flow not only a sampler, but a dynamical
tool for renormalization-group analysis.

\appendix

\section{Working Notes on Overlap, Locality, and Scalability}
\label{app:working_notes}

This appendix records the current working interpretation of the project and the
near-term exploratory program. These points are intended as guidance for the
full paper rather than as finalized claims.

\subsection{Relation to the Shorter Methods Paper}

There will inevitably be some overlap between the shorter methods paper and the
present RG paper, because the RG interpretation depends on the same Wilsonian
transport architecture. The intended division of labor is:
\begin{itemize}
\item the shorter paper gives the detailed model description, the learned
Gaussian-prior construction, and the fine-level diagnostics,
\item the present paper gives only a compact recap of the model and then focuses
on the exact blocked hierarchy of the learned action, observable-space RG
trajectories, locality, and universality.
\end{itemize}
In practice this means that paper~2 should contain one concise section
describing the blocking map and conditional fluctuation trivialization, while
referring the reader to the methods paper for implementation details and
architecture ablations.

\subsection{Locality Versus Ultra-Locality}

For RG arguments the relevant property is not ultra-locality but locality in the
lattice-field-theory sense, namely exponential decay of induced couplings or
responses with distance. Ultra-locality would mean strictly finite range. The
learned action generated by the present hierarchy is not expected to be
ultra-local. The right question is whether it is local, in the sense that its
effective couplings decay exponentially.

The current architecture has a built-in locality advantage: the blocking map is
fixed and local, and the conditional fluctuation updates at each level use only
local neighborhoods on the appropriate coarse lattice. However, this does not by
itself prove locality of the induced fine-level action $S_\theta$. That
property should therefore be checked numerically rather than assumed.

The cleanest proposed locality diagnostics are:
\begin{enumerate}
\item score-response locality, based on the decay with separation of
\begin{equation}
\frac{\delta}{\delta \phi_y}\frac{\partial S_\theta}{\partial \phi_x},
\end{equation}
\item Hessian locality, based on the distance dependence of
\begin{equation}
H_{xy}(\phi)=\frac{\partial^2 S_\theta}{\partial \phi_x\partial \phi_y},
\end{equation}
averaged on representative field samples,
\item stability of blocked observables under changes of the microscopic target
point and under replacement of the microscopic target theory by the learned
theory.
\end{enumerate}
Exponential decay in the first two diagnostics would provide direct evidence
that the learned theory is local even if it is not ultra-local.

\subsection{Why Two Dimensions Require Care}

The locality question is tied to universality. One would like to argue that the
additional operators generated by the learned action do not spoil the
long-distance physics because they are irrelevant. In two-dimensional
$\phi^4$ theory, however, naive Gaussian power counting is not the right guide
near the critical line. The critical infrared theory is in the Ising
universality class, and relevance or irrelevance should ultimately be judged in
that framework, especially in the $Z_2$-even sector relevant here.

For this reason the practical universality argument should not rely only on an
operator-by-operator classification of the microscopic learned action. A safer
strategy is:
\begin{itemize}
\item establish locality of the learned action in the sense of exponential decay,
\item preserve the exact $Z_2$ symmetry,
\item show that blocked trajectories in observable space approach the same
renormalized trajectory as the target theory.
\end{itemize}

\subsection{Scalability at Fixed Intensive Error}

The raw quantity $\std(\Delta S)$ is extensive enough that it should not be used
alone to compare different volumes. The present working hypothesis is that the
right intensive measure of quality is
\begin{equation}
\frac{\std(\Delta S)}{L}
\end{equation}
in two dimensions, or equivalently $\std(\Delta S)/\sqrt{V}$. Existing
numerical evidence already suggests that this quantity is much more stable under
volume scaling than $\std(\Delta S)$ itself.

For physically relevant scaling tests, however, the comparison should
ultimately be organized at approximately fixed $\xi/L$, not at fixed bare
couplings. The reason is that in a finite-size scaling regime the natural
continuum-oriented notion of ``same physics'' is to keep the ratio of the
correlation length to the box size approximately constant while the volume is
varied. The observed $L$-scaling of $\std(\Delta S)/L$ should therefore be
interpreted in that setting.

This point also suggests a possible connection to standard critical
phenomenology. Writing the mismatch as a sum of local terms,
\begin{equation}
\Delta S(\phi)=\sum_x \delta \ell_x(\phi),
\end{equation}
one has
\begin{equation}
\mathrm{Var}(\Delta S)=\sum_{x,y}
\langle \delta \ell_x \delta \ell_y\rangle_c
=
V\sum_r C_{\delta \ell}(r).
\end{equation}
Away from criticality, the integrated correlator $\sum_r C_{\delta \ell}(r)$ is
finite, so $\std(\Delta S)\propto L$ in two dimensions and
$\std(\Delta S)/L$ is genuinely intensive. Near criticality, however, the
integrated correlator can acquire a critical enhancement. For a $Z_2$-symmetric
ansatz, the dominant overlap of $\delta \ell$ should lie in the even sector and
is plausibly energy-like. In that case one expects the critical enhancement of
$\std(\Delta S)/L$ to be controlled by the same susceptibility that governs the
specific heat. For the two-dimensional Ising universality class, where
$\alpha=0$ and $\nu=1$, this suggests a logarithmic growth,
\begin{equation}
\frac{\std(\Delta S)}{L}\sim \sqrt{\log \xi},
\end{equation}
or, at criticality in finite volume, $\std(\Delta S)/L \sim \sqrt{\log L}$ when
$\xi/L$ is held fixed. This scaling scenario should be checked directly against
the HMC determination of $\xi$.

An analogous argument applies to the fluctuations of the training gradient.
Each parameter derivative of the per-configuration loss is itself extensive:
\begin{equation}
g_a(\phi)=\partial_{\theta_a}\mathcal L(\phi;\theta)=\sum_x O_a(x).
\end{equation}
Therefore
\begin{equation}
\mathrm{Var}(g_a)=V\sum_r C_a(r),
\end{equation}
with $C_a(r)=\langle O_a(0)O_a(r)\rangle_c$. Away from criticality,
$\std(g_a)\sim \sqrt{V}$ while the fluctuation of the batch-averaged gradient
density scales like $(BV)^{-1/2}$ for batch size $B$. Near criticality the same
integrated-correlator enhancement appears. This suggests that any increase of
the batch size required for stable optimization near the critical line should be
interpreted as a physical effect tied to long-distance fluctuations of the
target theory, not only as an optimizer artifact.

This observation points to a sharper algorithmic question: what is the training
cost required to reach a fixed value of $\std(\Delta S)/L$ as the volume grows?
If the answer is polynomial in $V$, then one has strong evidence that the
learned Wilsonian hierarchy avoids the practical exponential barrier associated
with direct fine-level reweighting. This would not yet be a formal proof that
critical slowing down is solved in full generality, but it would be compelling
evidence that the cost of constructing a local representative of the same
renormalized trajectory remains manageable.

The natural scaling experiment is therefore to pick one or more intensive target
thresholds, for example
\begin{equation}
\std(\Delta S)/L \in \{0.06,\,0.05,\,0.04\},
\end{equation}
and then measure the optimizer steps and wall-clock time required to reach these
thresholds at
\begin{equation}
L=16,\ 32,\ 64,\ 128,\ \ldots
\end{equation}
using the same staged training protocol.

\subsection{Long-Term Direction}

If the locality and fixed-intensive-error scaling tests both succeed, then the
long-term project becomes more ambitious than exact correction of a microscopic
target action. One may instead regard the learned hierarchy as a practical route
to constructing local representatives of a target universality class and then
study or simulate those representatives directly. In that sense the present
program would turn the original flow-based objective on its head: the goal would
no longer be only to correct the microscopic theory, but to operationally follow
the renormalized trajectory itself.

\bibliographystyle{apsrev4-2}
\bibliography{refs}

\end{document}
